% ch6.tex — 练习与补充 II（中文版）
\chapter{练习与补充 II}

\section{常微分方程与矩阵}

二阶 ODE $y''+\omega^2 y=\beta$ 映射为 $\dot{\vect z}=\hat M\vect z+\vect\phi$，其中
$\hat M=\mat{0&-\omega^2\\1&0}$，回到第 2 章的谐振子矩阵。通解为
\begin{equation}
\vect z(t) = e^{\hat M t}\vect z(0) + \int_0^t e^{\hat M(t-t')}\vect\phi\,dt'.
\label{eq:6.1.1}
\end{equation}

对于 $\vect\phi=\vect 0$ 且 $\hat M$ 为常数，
\begin{equation}
e^{\hat M t} =
\mat{\cos(\omega t)&-\omega\sin(\omega t)\\
     \frac{1}{\omega}\sin(\omega t)&\cos(\omega t)},
\label{eq:6.1.2}
\end{equation}
与第 2 章推导的结果完全相同。

\begin{exercise}
用矩阵方法求解 $y''+2\zeta\omega_0 y'+\omega_0^2 y = F_0\cos(\Omega t)$，并推导共振曲线。
\end{exercise}

\section{Crofton--Glaisher 恒等式与热型方程}

Crofton 恒等式
$e^{\xi\hat A}\hat B = \hat B e^{\xi\hat A} + \xi[\hat A,\hat B]e^{\xi\hat A}$
应用于热多项式给出 Glaisher 恒等式
\begin{equation}
e^{at\partial_x^2}\,e^{-x^2}
= \frac{1}{\sqrt{1+4at}}\;
\exp\!\left(-\frac{x^2}{1+4at}\right).
\label{eq:6.2.1}
\end{equation}

一个推广：对于热型方程
$\partial_t F = a\,\partial_x^2 F$ 且 $F(x,0)=x^n$，
\begin{equation}
F(x,t) = H_n(x,at) = n!\sum_{r=0}^{\lfloor n/2\rfloor}
\frac{x^{n-2r}(at)^r}{(n-2r)!\,r!}.
\label{eq:6.2.2}
\end{equation}

\begin{exercise}
证明 Crofton 恒等式 $e^{y\partial_x^m}f(x)=f(x+my\partial_x^{m-1})e^{y\partial_x^m}$，
并利用它推导 $\partial_x^{2m}$ 情形的广义 Glaisher 恒等式。
\end{exercise}

\section{Gamma 函数与定积分}

Euler 的 Gamma 函数：
\begin{equation}
\Gamma(z) = \int_0^\infty t^{z-1}e^{-t}\,dt, \qquad \Re(z)>0.
\label{eq:6.3.1}
\end{equation}

关键性质：
\begin{align}
\Gamma(z+1) &= z\,\Gamma(z),\\
\Gamma(\tfrac12) &= \sqrt{\pi},\\
\Gamma(n+1) &= n! \quad (n\in\mathbb N_0).
\label{eq:6.3.2}
\end{align}

Beta 函数：
\begin{equation}
B(p,q) = \int_0^1 t^{p-1}(1-t)^{q-1}dt
= \frac{\Gamma(p)\Gamma(q)}{\Gamma(p+q)}.
\label{eq:6.3.3}
\end{equation}

有用的积分：
\begin{equation}
\int_0^\infty x^{n}e^{-ax^2}dx =
\frac{\Gamma\!\left(\frac{n+1}{2}\right)}{2a^{(n+1)/2}},
\qquad a>0,\;n>-1.
\label{eq:6.3.4}
\end{equation}

\begin{exercise}
用 Beta 函数计算 $\int_0^{\pi/2}\sin^{2m}\theta\cos^{2n}\theta\,d\theta$。
\end{exercise}

\section{复变函数方法与积分计算}

对于在简单区域中解析的函数 $f(z)$，闭合路径积分消失（Cauchy 定理）。留数定理：
\begin{equation}
\oint_C f(z)\,dz = 2\pi i\sum_{z_k\in\text{int}(C)} \operatorname{Res}(f,z_k).
\label{eq:6.4.1}
\end{equation}

Jordan 引理：若 $f(z)$ 在上/下半平面内当 $|z|\to\infty$ 时一致趋于 $0$，则半圆弧积分消失。

\begin{exercise}
利用上半平面中的半圆围道计算 $\int_{-\infty}^{\infty}\frac{\cos(kx)}{1+x^2}\,dx$。
\end{exercise}

主值：对于实轴上 $x=\Omega$ 处的极点，
\begin{equation}
\mathcal P\int_{-\infty}^{\infty}\frac{\tilde\chi(\omega)}{\omega-\Omega}\,d\omega
= \lim_{\epsilon\to0}\left(
\int_{-\infty}^{\Omega-\epsilon}+\int_{\Omega+\epsilon}^{\infty}
\right)\frac{\tilde\chi(\omega)}{\omega-\Omega}\,d\omega.
\label{eq:6.4.2}
\end{equation}
$\omega=\Omega$ 周围的小半圆贡献为 $-i\pi\tilde\chi(\Omega)$。

\section{Fourier 变换}

定义：
\begin{equation}
\tilde f(k)=\int_{-\infty}^{\infty}f(x)\,e^{-ikx}\,dx,\qquad
f(x)=\frac{1}{2\pi}\int_{-\infty}^{\infty}\tilde f(k)\,e^{ikx}\,dk.
\label{eq:6.5.1}
\end{equation}

该变换把 $\partial_x$ 对角化为 $ik$，并把 PDE 化为 ODE。关键性质：
\begin{align}
\mathcal F[f'(x)] &= ik\,\tilde f(k),\\
\mathcal F[xf(x)] &= i\,\widetilde f\,'(k),\\
\mathcal F[f*g] &= \tilde f(k)\,\tilde g(k),
\label{eq:6.5.2}
\end{align}
其中 $(f*g)(x)=\int f(y)g(x-y)\,dy$。

Gauss 函数：
\begin{equation}
\mathcal F[e^{-ax^2}] = \sqrt{\frac{\pi}{a}}\,e^{-k^2/(4a)}.
\label{eq:6.5.3}
\end{equation}

\begin{exercise}
推导卷积定理，并用它通过 Fourier 变换求解
$y''+k_0^2y=s(x)$。
\end{exercise}

\section{Fourier 变换与微分方程的求解}

对于边界条件在无穷远的 $y''+k_0^2y=s(x)$，
\begin{equation}
\tilde y(k) = \frac{\tilde s(k)}{k_0^2-k^2},
\label{eq:6.6.1}
\end{equation}
逆变换给出 Green 函数解
\begin{equation}
y(x) = \int_{-\infty}^{\infty} G(x-\xi)\,s(\xi)\,d\xi,
\qquad G(x) = \frac{1}{2\pi}\int_{-\infty}^{\infty}
\frac{e^{ikx}}{k_0^2-k^2}\,dk.
\label{eq:6.6.2}
\end{equation}

在上/下半平面闭合围道给出
$G(x)=\frac{\sin(k_0|x|)}{k_0}$，回到第 8 章的 Green 函数。

\section{Fourier 型变换}

Hankel 变换（柱对称）：
\begin{equation}
\tilde f(\rho) = \int_0^\infty f(r)\,J_0(\rho r)\,r\,dr,
\qquad
f(r) = \int_0^\infty \tilde f(\rho)\,J_0(\rho r)\,\rho\,d\rho.
\label{eq:6.7.1}
\end{equation}

Mellin 变换：
\begin{equation}
\mathcal M[f](s) = \int_0^\infty f(x)\,x^{s-1}\,dx,
\qquad
f(x) = \frac{1}{2\pi i}\int_{c-i\infty}^{c+i\infty}
\mathcal M[f](s)\,x^{-s}\,ds.
\label{eq:6.7.2}
\end{equation}

Laplace 变换：
\begin{equation}
\mathcal L[f](s) = \int_0^\infty f(t)\,e^{-st}\,dt,
\qquad
\mathcal L[f'](s) = s\tilde f(s)-f(0).
\label{eq:6.7.3}
\end{equation}

\begin{exercise}
证明 Hankel 变换是柱坐标中的 Fourier 变换，并推导 $J_0$ 的正交关系。
\end{exercise}

% 图形
\begin{figure}[h]
\centering
\includegraphics[max width=0.85\linewidth,keepaspectratio]{fig6_1.png}
\caption{Fourier 变换对：Gauss 函数及其变换。}
\label{fig:6.1}
\end{figure}

\begin{figure}[h]
\centering
\includegraphics[max width=0.85\linewidth,keepaspectratio]{fig6_13.png}
\caption{主值积分在复 $\omega$ 平面中的积分围道。}
\label{fig:6.13}
\end{figure}
